\documentclass[border=8pt]{standalone}
% c-comp/latex/figure-preamble.tex の配色と LuaLaTeX + TikZ の方式を参照。
\usepackage{luatexja}
\usepackage{luatexja-fontspec}
\setmainfont{Noto Sans CJK JP}
\setmainjfont{Noto Sans CJK JP}
\usepackage{amsmath,amssymb}
\usepackage{tikz}
\usetikzlibrary{arrows.meta,calc,patterns,positioning}
\definecolor{Accent}{HTML}{583B66}
\definecolor{Light}{HTML}{EEE6F2}
\definecolor{Emph}{HTML}{C44919}
\definecolor{Ink}{HTML}{262626}
\tikzset{
  every node/.style={text=Ink,font=\small},
  heading/.style={font=\small\bfseries,text=Accent,align=center},
  cell/.style={draw=Accent,minimum width=7mm,minimum height=7mm,inner sep=0pt},
  flow/.style={-{Stealth[length=2mm]},draw=Accent,thick},
  note/.style={align=left,font=\small},
}

\begin{document}
\begin{tikzpicture}[x=7mm,y=7mm]
  \node[heading] at (2,1.8) {入力 $a$：勤務可能性};
  \node[heading] at (8,1.8) {勤務表 $x$：解の一例};
  \foreach \offset in {0,6} {
    \foreach \d in {1,2,3} {
      \node at (\offset+\d,0.85) {$\d$};
    }
    \foreach \p in {1,2,3,4} {
      \node[anchor=east] at (\offset+0.35,1-\p) {$p=\p$};
      \foreach \d in {1,2,3} {
        \node[cell] at (\offset+\d,1-\p) {};
      }
    }
  }
  % 入力0と、対応する解の0は斜線。入力1でも解0は許される。
  \foreach \p/\d in {1/3,2/2,3/1} {
    \foreach \offset in {0,6} {
      \path[pattern=north east lines,pattern color=Light]
        (\offset+\d-.5,1-\p-.5) rectangle (\offset+\d+.5,1-\p+.5);
    }
  }
  \foreach \p/\d/\v in {1/1/1,1/2/1,1/3/0,2/1/1,2/2/0,2/3/1,3/1/0,3/2/1,3/3/1,4/1/1,4/2/1,4/3/1} {
    \node at (\d,1-\p) {$\v$};
  }
  \foreach \p/\d/\v in {1/1/1,1/2/1,1/3/0,2/1/1,2/2/0,2/3/1,3/1/0,3/2/1,3/3/1,4/1/0,4/2/0,4/3/0} {
    \ifnum\v=1
      \node[cell,fill=Light] at (6+\d,1-\p) {$\v$};
    \else
      \node at (6+\d,1-\p) {$\v$};
    \fi
  }
  \draw[Emph,very thick] (0.5,-3.5) rectangle (1.5,-2.5);
  \draw[Emph,very thick] (6.5,-3.5) rectangle (7.5,-2.5);
  \node at (4.5,-4) {行：人 $p$ ／ 列：日 $d$};
  \node[note,anchor=north west,text width=80mm] at (0,-4.7)
    {斜線のマス：勤務不可なので、勤務表でも必ず0。\\
     塗ったマス：実際に勤務する1。\\
     強調枠 $(p,d)=(4,1)$：勤務可能な1でも、選ばず0でよい。};
  \node[heading,anchor=north] at (4.5,-7.5)
    {$x_{p,d}\le a_{p,d}$：入力と同じ値にする条件ではない};
\end{tikzpicture}
\end{document}
